% Options for packages loaded elsewhere
% Options for packages loaded elsewhere
\PassOptionsToPackage{unicode,draft}{hyperref}
\PassOptionsToPackage{hyphens}{url}
\PassOptionsToPackage{dvipsnames,svgnames,x11names}{xcolor}
%
\documentclass[
  letterpaper,
  DIV=11,
  numbers=noendperiod]{scrartcl}
\usepackage{xcolor}
\usepackage{amsmath,amssymb}
\setcounter{secnumdepth}{-\maxdimen} % remove section numbering
\usepackage{iftex}
\ifPDFTeX
  \usepackage[T1]{fontenc}
  \usepackage[utf8]{inputenc}
  \usepackage{textcomp} % provide euro and other symbols
\else % if luatex or xetex
  \usepackage{unicode-math} % this also loads fontspec
  \defaultfontfeatures{Scale=MatchLowercase}
  \defaultfontfeatures[\rmfamily]{Ligatures=TeX,Scale=1}
\fi
\usepackage{lmodern}
\ifPDFTeX\else
  % xetex/luatex font selection
\fi
% Use upquote if available, for straight quotes in verbatim environments
\IfFileExists{upquote.sty}{\usepackage{upquote}}{}
\IfFileExists{microtype.sty}{% use microtype if available
  \usepackage[]{microtype}
  \UseMicrotypeSet[protrusion]{basicmath} % disable protrusion for tt fonts
}{}
\makeatletter
\@ifundefined{KOMAClassName}{% if non-KOMA class
  \IfFileExists{parskip.sty}{%
    \usepackage{parskip}
  }{% else
    \setlength{\parindent}{0pt}
    \setlength{\parskip}{6pt plus 2pt minus 1pt}}
}{% if KOMA class
  \KOMAoptions{parskip=half}}
\makeatother
% Make \paragraph and \subparagraph free-standing
\makeatletter
\ifx\paragraph\undefined\else
  \let\oldparagraph\paragraph
  \renewcommand{\paragraph}{
    \@ifstar
      \xxxParagraphStar
      \xxxParagraphNoStar
  }
  \newcommand{\xxxParagraphStar}[1]{\oldparagraph*{#1}\mbox{}}
  \newcommand{\xxxParagraphNoStar}[1]{\oldparagraph{#1}\mbox{}}
\fi
\ifx\subparagraph\undefined\else
  \let\oldsubparagraph\subparagraph
  \renewcommand{\subparagraph}{
    \@ifstar
      \xxxSubParagraphStar
      \xxxSubParagraphNoStar
  }
  \newcommand{\xxxSubParagraphStar}[1]{\oldsubparagraph*{#1}\mbox{}}
  \newcommand{\xxxSubParagraphNoStar}[1]{\oldsubparagraph{#1}\mbox{}}
\fi
\makeatother


\usepackage{longtable,booktabs,array}
\usepackage{calc} % for calculating minipage widths
% Correct order of tables after \paragraph or \subparagraph
\usepackage{etoolbox}
\makeatletter
\patchcmd\longtable{\par}{\if@noskipsec\mbox{}\fi\par}{}{}
\makeatother
% Allow footnotes in longtable head/foot
\IfFileExists{footnotehyper.sty}{\usepackage{footnotehyper}}{\usepackage{footnote}}
\makesavenoteenv{longtable}
\usepackage{graphicx}
\makeatletter
\newsavebox\pandoc@box
\newcommand*\pandocbounded[1]{% scales image to fit in text height/width
  \sbox\pandoc@box{#1}%
  \Gscale@div\@tempa{\textheight}{\dimexpr\ht\pandoc@box+\dp\pandoc@box\relax}%
  \Gscale@div\@tempb{\linewidth}{\wd\pandoc@box}%
  \ifdim\@tempb\p@<\@tempa\p@\let\@tempa\@tempb\fi% select the smaller of both
  \ifdim\@tempa\p@<\p@\scalebox{\@tempa}{\usebox\pandoc@box}%
  \else\usebox{\pandoc@box}%
  \fi%
}
% Set default figure placement to htbp
\def\fps@figure{htbp}
\makeatother





\setlength{\emergencystretch}{3em} % prevent overfull lines

\providecommand{\tightlist}{%
  \setlength{\itemsep}{0pt}\setlength{\parskip}{0pt}}



 


\KOMAoption{captions}{tableheading}
\usepackage{fancyvrb}
\usepackage{fvextra}
\RecustomVerbatimEnvironment{verbatim}{Verbatim}{breaklines=true,breakanywhere=true,fontsize=\small}
\makeatletter
\@ifpackageloaded{caption}{}{\usepackage{caption}}
\AtBeginDocument{%
\ifdefined\contentsname
  \renewcommand*\contentsname{Table of contents}
\else
  \newcommand\contentsname{Table of contents}
\fi
\ifdefined\listfigurename
  \renewcommand*\listfigurename{List of Figures}
\else
  \newcommand\listfigurename{List of Figures}
\fi
\ifdefined\listtablename
  \renewcommand*\listtablename{List of Tables}
\else
  \newcommand\listtablename{List of Tables}
\fi
\ifdefined\figurename
  \renewcommand*\figurename{Figure}
\else
  \newcommand\figurename{Figure}
\fi
\ifdefined\tablename
  \renewcommand*\tablename{Table}
\else
  \newcommand\tablename{Table}
\fi
}
\@ifpackageloaded{float}{}{\usepackage{float}}
\floatstyle{ruled}
\@ifundefined{c@chapter}{\newfloat{codelisting}{h}{lop}}{\newfloat{codelisting}{h}{lop}[chapter]}
\floatname{codelisting}{Listing}
\newcommand*\listoflistings{\listof{codelisting}{List of Listings}}
\makeatother
\makeatletter
\makeatother
\makeatletter
\@ifpackageloaded{caption}{}{\usepackage{caption}}
\@ifpackageloaded{subcaption}{}{\usepackage{subcaption}}
\makeatother
\usepackage{bookmark}
\IfFileExists{xurl.sty}{\usepackage{xurl}}{} % add URL line breaks if available
\urlstyle{same}
\hypersetup{
  pdftitle={samling},
  pdfauthor={st. schwarz},
  colorlinks=true,
  linkcolor={blue},
  filecolor={Maroon},
  citecolor={Blue},
  urlcolor={Blue},
  pdfcreator={LaTeX via pandoc}}


\title{samling}
\usepackage{etoolbox}
\makeatletter
\providecommand{\subtitle}[1]{% add subtitle to \maketitle
  \apptocmd{\@title}{\par {\large #1 \par}}{}{}
}
\makeatother
\subtitle{draft B samling}
\author{st. schwarz}
\date{2026-02-28}
\begin{document}
\maketitle


\section{index}\label{index}

\section{GPT essais}\label{gpt-essais}

\section{bezug}\label{bezug}

\begin{figure}[H]

{\centering \pandocbounded{\includegraphics[keepaspectratio]{play_files/mediabag/gpt-001-endetag.jpg}}

}

\caption{Q: transformer generated image}

\end{figure}%

\section{prompt 1:}\label{prompt-1}

``frage: wird tim\&struppi oft als \{mask\} für solche illustrationen
verwendet? kannst du das fehlende wort im satz ausfüllen und die frage
dann beantworten?''

\subsection{response}\label{response}

Das fehlende Wort, das hier \textbf{am besten passt}, ist:

\begin{quote}
\textbf{Projektionsfläche}
\end{quote}

\textbf{Vollständige Frage:}

\begin{quote}
\emph{Wird Tim \& Struppi oft als \textbf{Projektionsfläche} für solche
Illustrationen verwendet?}
\end{quote}

\subsubsection{Antwort:}\label{antwort}

\textbf{Ja -- sehr häufig.}

Tim \& Struppi (insbesondere \textbf{Tim allein, ohne Struppi}) wird
seit Jahrzehnten als Projektionsfläche genutzt, weil:

\begin{itemize}
\tightlist
\item
  Tim \textbf{kaum individuelle Eigenschaften} hat → keine ausgeprägte
  Psychologie, keine Entwicklung, kaum Ironie
\item
  er bewusst als \textbf{neutraler Beobachter} angelegt ist
\item
  sein Design extrem \textbf{reduziert und archetypisch} ist
\item
  er leicht in andere Rollen „übersetzt`` werden kann (Reporter,
  Abenteurer, Cowboy, Reisender, Stellvertreter des Lesers)
\end{itemize}

Dadurch eignet sich die Figur (oder eine Anspielung darauf) perfekt für:

\begin{itemize}
\tightlist
\item
  Parodien
\item
  politische Karikaturen
\item
  philosophische Pointen
\item
  absurde oder meta-hafte Aussagen
\item
  „unschuldige Figur in widersprüchlichem Kontext``
\end{itemize}

Genau das machst du hier:

\begin{itemize}
\tightlist
\item
  eine \textbf{naiv wirkende Comicfigur}
\item
  in einer \textbf{moralisch / logisch widersprüchlichen Gedankenblase}
\item
  → Reibung erzeugt Humor
\end{itemize}

⚠️ \textbf{Rechtlich} ist genau diese Stärke zugleich die Schwäche:

\begin{itemize}
\tightlist
\item
  Weil Tim so klar erkennbar ist, reagieren Rechteinhaber sensibel.
\item
  Deshalb ist die Strategie „Anspielung statt Figur`` (wie in deiner
  letzten Version) genau richtig.
\end{itemize}

Wenn du magst, kann ich dir auch:

\begin{itemize}
\tightlist
\item
  2--3 \textbf{alternative Archetypen} nennen, die ähnlich
  funktionieren,
\item
  aber \textbf{noch weniger rechtliches Risiko} haben.
\end{itemize}

\section{prompt 2}\label{prompt-2}

noch: kannst du den teil ``am ende des tages'' in der blase kursiv
setzen und dann hier auf den unterschied der letzten beiden versionen
hinsichtlich linguistischer marker eingehen, also was die jeweiligen
versionen linguistisch unterscheidet.

\subsection{response}\label{response-1}

Gerne. Ich mache das in zwei klar getrennten Schritten.

\begin{center}\rule{0.5\linewidth}{0.5pt}\end{center}

\subsection{1. Denkblase mit kursivem
Marker}\label{denkblase-mit-kursivem-marker}

\begin{quote}
\emph{Am Ende des Tages} sollte der Hund auch gegessen werden.
\end{quote}

Die Kursivsetzung markiert \textbf{explizit} den idiomatischen Teil der
Aussage.

\begin{center}\rule{0.5\linewidth}{0.5pt}\end{center}

\subsection{2. Linguistischer Unterschied der letzten beiden
Versionen}\label{linguistischer-unterschied-der-letzten-beiden-versionen}

(„Hund am Spieß`` → „Gemüse am Spieß``)

Der Unterschied liegt \textbf{nicht} primär im Bildinhalt, sondern in
den \textbf{linguistischen Markern}, die das Bild \emph{lesen} lassen.

\subsubsection{Version A (Hund / gegrilltes
Tier)}\label{version-a-hund-gegrilltes-tier}

\textbf{Linguistische Wirkung: direkte Semantik}

\begin{itemize}
\item
  „Hund`` wird \textbf{wörtlich} gelesen
\item
  Das Idiom „am Ende des Tages`` verliert seine Metaphernfunktion
\item
  Die Aussage wird als \textbf{propositionale Behauptung} interpretiert:

  \begin{quote}
  X ist ein Hund → X wird gegessen
  \end{quote}
\end{itemize}

➡️ Linguistisch dominant:

\begin{itemize}
\tightlist
\item
  \textbf{Denotation}
\item
  geringe Distanz zwischen Wort und Referent
\item
  kaum Ironiesignal
\item
  Schock / Tabubruch entsteht \textbf{außersprachlich}
\end{itemize}

\begin{center}\rule{0.5\linewidth}{0.5pt}\end{center}

\subsubsection{Version B (Gemüse /
Parodie)}\label{version-b-gemuxfcse-parodie}

\textbf{Linguistische Wirkung: idiomatische Reibung}

Hier greifen mehrere Marker gleichzeitig:

\paragraph{\texorpdfstring{a) \emph{Am Ende des Tages}
(kursiv)}{a) Am Ende des Tages (kursiv)}}\label{a-am-ende-des-tages-kursiv}

\begin{itemize}
\item
  expliziter \textbf{Diskursmarker}
\item
  signalisiert:

  \begin{itemize}
  \tightlist
  \item
    Zusammenfassung
  \item
    Relativierung
  \item
    rhetorische Routine
  \end{itemize}
\item
  wird als \textbf{formelhafte Leerformel} erkennbar
\end{itemize}

\paragraph{b) „Hund`` ohne Hund}\label{b-hund-ohne-hund}

\begin{itemize}
\item
  Referent fehlt oder ist verschoben
\item
  „Hund`` wird \textbf{semantisch entkoppelt}
\item
  Leser muss interpretieren:

  \begin{itemize}
  \tightlist
  \item
    Metapher?
  \item
    Redensart?
  \item
    absurde Übertragung?
  \end{itemize}
\end{itemize}

\paragraph{c) Bild--Text-Divergenz}\label{c-bildtext-divergenz}

\begin{itemize}
\tightlist
\item
  Bild zeigt harmloses, vegetarisches Grillen
\item
  Text evoziert etwas völlig anderes
\end{itemize}

➡️ Linguistisch dominant:

\begin{itemize}
\tightlist
\item
  \textbf{Pragmatik statt Semantik}
\item
  Ironie durch \textbf{Inkongruenz}
\item
  Leser erkennt: \emph{Das ist nicht wörtlich gemeint}
\end{itemize}

\begin{center}\rule{0.5\linewidth}{0.5pt}\end{center}

\subsection{3. Kernunterschied in einem
Satz}\label{kernunterschied-in-einem-satz}

\begin{itemize}
\tightlist
\item
  \textbf{Version A} zwingt eine \textbf{wörtliche Lesart} auf.
\item
  \textbf{Version B} erzwingt eine \textbf{metasprachliche Lesart}.
\end{itemize}

Oder noch präziser:

\begin{quote}
In Version B wird der Satz nicht mehr \emph{verstanden}, sondern
\emph{interpretiert}.
\end{quote}

\begin{center}\rule{0.5\linewidth}{0.5pt}\end{center}

\subsection{4. Warum die Kursivsetzung entscheidend
ist}\label{warum-die-kursivsetzung-entscheidend-ist}

Die Kursivierung von \emph{Am Ende des Tages} wirkt wie ein Schild:

\begin{quote}
„Achtung: Das Folgende ist Rhetorik, nicht Realität.``
\end{quote}

Sie macht sichtbar, dass:

\begin{itemize}
\tightlist
\item
  hier eine \textbf{sprachliche Konvention} zitiert wird
\item
  nicht eine Handlung beschrieben werden soll
\end{itemize}

Damit verschiebst du die Aussage:

\begin{itemize}
\tightlist
\item
  von \textbf{Ethik → Sprachkritik}
\item
  von \textbf{Provokation → Reflexion}
\end{itemize}

\begin{center}\rule{0.5\linewidth}{0.5pt}\end{center}

\subsubsection{Kurzfazit}\label{kurzfazit}

Die letzte Version ist linguistisch:

\begin{itemize}
\tightlist
\item
  ironischer
\item
  reflektierter
\item
  klarer als Parodie lesbar
\item
  und deutlich robuster gegenüber Missverständnissen
\end{itemize}

Wenn du willst, können wir den Satz auch noch \textbf{minimal
umstellen}, um die Ironie weiter zu schärfen, ohne den Wortlaut stark zu
verändern.

\section{prompt 3}\label{prompt-3}

wie würdest du das machen und auf welche trainierten ressourcen greifst
du bei der erkennung/generierung der ironie zu?

\subsection{response}\label{response-2}

Gute, anspruchsvolle Frage. Ich teile die Antwort in \textbf{(A)
konkrete sprachliche Eingriffe} und \textbf{(B) Meta-Ebene: wie ich
Ironie erkenne und erzeuge}.

\begin{center}\rule{0.5\linewidth}{0.5pt}\end{center}

\subsection{A) Wie ich den Satz minimal verändern würde (ohne die Pointe
zu
verlieren)}\label{a-wie-ich-den-satz-minimal-veruxe4ndern-wuxfcrde-ohne-die-pointe-zu-verlieren}

Ausgangspunkt:

\begin{quote}
\emph{Am Ende des Tages} sollte der Hund auch gegessen werden.
\end{quote}

\subsubsection{1. Minimalverschiebung durch
Modalabschwächung}\label{minimalverschiebung-durch-modalabschwuxe4chung}

\begin{quote}
\emph{Am Ende des Tages} \textbf{muss} der Hund ja auch gegessen werden.
\end{quote}

\textbf{Effekt:}

\begin{itemize}
\tightlist
\item
  „ja auch`` ist ein klassischer \textbf{Ironiemarker}
\item
  signalisiert angebliche Selbstverständlichkeit
\item
  erzeugt Distanz zwischen Sprecher und Aussage
\end{itemize}

\begin{center}\rule{0.5\linewidth}{0.5pt}\end{center}

\subsubsection{2. Parenthetische
Brechung}\label{parenthetische-brechung}

\begin{quote}
\emph{Am Ende des Tages} sollte der Hund \textbf{wohl} auch gegessen
werden.
\end{quote}

\textbf{Effekt:}

\begin{itemize}
\tightlist
\item
  „wohl`` markiert epistemische Unsicherheit
\item
  unterminiert die Aussage von innen heraus
\item
  typisch für ironische Selbstrelativierung
\end{itemize}

\begin{center}\rule{0.5\linewidth}{0.5pt}\end{center}

\subsubsection{3. Bürokratische Überformung (sehr
stark)}\label{buxfcrokratische-uxfcberformung-sehr-stark}

\begin{quote}
\emph{Am Ende des Tages} ist der Hund auch \textbf{zu verwerten}.
\end{quote}

\textbf{Effekt:}

\begin{itemize}
\tightlist
\item
  technokratische Sprache
\item
  moralische Kälte durch Nominalisierung
\item
  klassische Satiretechnik (→ Kafka, Loriot, Verwaltungsdeutsch)
\end{itemize}

\begin{center}\rule{0.5\linewidth}{0.5pt}\end{center}

\subsubsection{4. Meta-Diskurs explizit machen (feiner
Humor)}\label{meta-diskurs-explizit-machen-feiner-humor}

\begin{quote}
\emph{Am Ende des Tages} sagt man ja, der Hund sollte auch gegessen
werden.
\end{quote}

\textbf{Effekt:}

\begin{itemize}
\tightlist
\item
  Sprecher schiebt Verantwortung auf ein anonymes „man``
\item
  markiert \textbf{Zitat statt Überzeugung}
\end{itemize}

\begin{center}\rule{0.5\linewidth}{0.5pt}\end{center}

\subsubsection{Meine Empfehlung für dein
Bild}\label{meine-empfehlung-fuxfcr-dein-bild}

\begin{quote}
\emph{Am Ende des Tages} muss der Hund ja auch gegessen werden.
\end{quote}

Die Pointe bleibt, aber die \textbf{Ironie ist linguistisch
abgesichert}.

\begin{center}\rule{0.5\linewidth}{0.5pt}\end{center}

\subsection{B) Auf welche „trainierten Ressourcen`` ich bei Ironie
zurückgreife}\label{b-auf-welche-trainierten-ressourcen-ich-bei-ironie-zuruxfcckgreife}

Ich greife \textbf{nicht} auf einzelne gespeicherte Beispiele zu,
sondern auf \textbf{verallgemeinerte Muster}, die aus sehr vielen Texten
gelernt wurden. Vereinfacht lassen sie sich in vier Ebenen einteilen:

\begin{center}\rule{0.5\linewidth}{0.5pt}\end{center}

\subsubsection{1. Idiom- und
Formelwissen}\label{idiom--und-formelwissen}

Ich erkenne:

\begin{itemize}
\item
  „am Ende des Tages`` = \textbf{diskursives Füllidiom}
\item
  hohe Frequenz in:

  \begin{itemize}
  \tightlist
  \item
    Managementsprache
  \item
    politischen Statements
  \item
    Rechtfertigungsrhetorik
  \end{itemize}
\end{itemize}

➡️ Hohe Wahrscheinlichkeit für \textbf{ironische Verwendung}, wenn
Kontext nicht passt.

\begin{center}\rule{0.5\linewidth}{0.5pt}\end{center}

\subsubsection{2. Inkonkurrenz-Detektion
(Kernmechanismus)}\label{inkonkurrenz-detektion-kernmechanismus}

Ironie entsteht fast immer durch \textbf{Mismatch} zwischen:

\begin{longtable}[]{@{}ll@{}}
\toprule\noalign{}
Ebene & Beispiel \\
\midrule\noalign{}
\endhead
\bottomrule\noalign{}
\endlastfoot
sprachlich & rationale Floskel \\
visuell & harmloses Grillen \\
moralisch & problematische Aussage \\
\end{longtable}

Ich bin darauf trainiert, solche \textbf{Reibungen} zu erkennen, weil
sie in Satire, Feuilleton, Literatur und Dialogen extrem häufig sind.

\begin{center}\rule{0.5\linewidth}{0.5pt}\end{center}

\subsubsection{3. Pragmatik statt Logik}\label{pragmatik-statt-logik}

Ironie ist selten logisch falsch, sondern \textbf{kommunikativ falsch}.

Ich bewerte:

\begin{itemize}
\tightlist
\item
  Wer spricht?
\item
  In welchem Ton?
\item
  Mit welcher erwartbaren Intention?
\item
  Welche Lesart wäre sozial „unpassend``?
\end{itemize}

Wenn eine Lesart \emph{zu glatt} ist, wird sie verdächtig.

\begin{center}\rule{0.5\linewidth}{0.5pt}\end{center}

\subsubsection{4. Marker-Lexikon (implizit
gelernt)}\label{marker-lexikon-implizit-gelernt}

Bestimmte Wörter fungieren statistisch als Ironiesignale:

\begin{itemize}
\tightlist
\item
  „ja``
\item
  „wohl``
\item
  „eben``
\item
  „natürlich``
\item
  „man sagt``
\end{itemize}

Diese Marker habe ich nicht als Liste, sondern als
\textbf{Wahrscheinlichkeitsverteilungen} gelernt:

\begin{quote}
Wenn X + Kontext Y → Ironie wahrscheinlicher als Literalität.
\end{quote}

\begin{center}\rule{0.5\linewidth}{0.5pt}\end{center}

\subsection{\texorpdfstring{C) Wichtig: Was ich \emph{nicht}
habe}{C) Wichtig: Was ich nicht habe}}\label{c-wichtig-was-ich-nicht-habe}

\begin{itemize}
\tightlist
\item
  kein eigenes Welturteil
\item
  kein echtes „Verstehen``
\item
  kein Zugriff auf einzelne Trainingsdokumente
\item
  keine Intuition im menschlichen Sinn
\end{itemize}

Was ich habe:

\begin{itemize}
\tightlist
\item
  extrem dichte \textbf{Korrelationen zwischen Form, Kontext und
  Wirkung}
\end{itemize}

\begin{center}\rule{0.5\linewidth}{0.5pt}\end{center}

\subsection{D) Ein Satz zur Essenz}\label{d-ein-satz-zur-essenz}

\begin{quote}
Ironie erkenne und erzeuge ich nicht, weil ich sie „verstehe``, sondern
weil ich gelernt habe, \textbf{wo Sprache typischerweise aufhört, ernst
gemeint zu sein}.
\end{quote}

Wenn du willst, können wir als Nächstes:

\begin{itemize}
\tightlist
\item
  die Denkblase noch feiner justieren \textbf{(ein Wort austauschen,
  große Wirkung)}
\item
  oder das Bild--Text-Verhältnis weiter zuspitzen (noch subtiler, noch
  sicherer).
\end{itemize}

\section{prompt 4}\label{prompt-4}

wie bewertest/eruierst du in punkt c) ``wirkung''?

\subsection{response}\label{response-3}

Sehr gute Nachfrage. „Wirkung`` ist der heikelste Punkt -- und zugleich
der interessanteste. Ich beantworte das \textbf{präzise, nicht
mystifizierend}.

\begin{center}\rule{0.5\linewidth}{0.5pt}\end{center}

\subsection{Kurzantwort}\label{kurzantwort}

Ich \textbf{messe} Wirkung nicht, ich \textbf{schätze sie
probabilistisch} -- als \emph{erwartbare Reaktion typischer Rezipienten}
in vergleichbaren Diskursen.

\begin{center}\rule{0.5\linewidth}{0.5pt}\end{center}

\subsection{Langfassung: Wie ich „Wirkung`` in Punkt C euruieren
kann}\label{langfassung-wie-ich-wirkung-in-punkt-c-euruieren-kann}

\subsubsection{1. Wirkung ≠ Reaktion}\label{wirkung-reaktion}

Zuerst eine Trennung:

\begin{itemize}
\tightlist
\item
  \textbf{Reaktion}: was ein konkreter Mensch tatsächlich fühlt
\item
  \textbf{Wirkung}: was ein sprachlicher Akt \emph{typischerweise}
  auslöst
\end{itemize}

Ich arbeite ausschließlich mit \textbf{Wirkung} im zweiten Sinn.

\begin{center}\rule{0.5\linewidth}{0.5pt}\end{center}

\subsection{2. Wirkung als
Erwartungswert}\label{wirkung-als-erwartungswert}

Formal gedacht (ohne Mathematik):

\begin{quote}
Wirkung = \textbf{(Sprachform × Kontext × kulturelle Konventionen)} →
\textbf{wahrscheinliche Interpretationsmodi}
\end{quote}

Ich schätze:

\begin{itemize}
\tightlist
\item
  Wie viele Lesarten sind wahrscheinlich?
\item
  Welche davon ist sozial dominant?
\item
  Welche ist markiert (auffällig, störend, komisch)?
\end{itemize}

\begin{center}\rule{0.5\linewidth}{0.5pt}\end{center}

\subsection{3. Die drei Ebenen der
Wirkungsabschätzung}\label{die-drei-ebenen-der-wirkungsabschuxe4tzung}

\subsubsection{Ebene A: Sprachinterne
Signale}\label{ebene-a-sprachinterne-signale}

Ich analysiere:

\begin{itemize}
\tightlist
\item
  Modalpartikeln („ja``, „wohl``, „eben``)
\item
  Idiome („am Ende des Tages``)
\item
  Tempus / Modus
\item
  Nominalisierungen
\end{itemize}

Beispiel:

\begin{quote}
„muss ja auch``
\end{quote}

➡️ Hohe Korrelation mit:

\begin{itemize}
\tightlist
\item
  Rechtfertigung
\item
  Nachträglicher Rationalisierung
\item
  ironischer Distanz
\end{itemize}

\textbf{Erwartete Wirkung:} Nicht Zustimmung, sondern Stirnrunzeln /
Lachen / Distanz.

\begin{center}\rule{0.5\linewidth}{0.5pt}\end{center}

\subsubsection{Ebene B: Kontextuelle
Plausibilität}\label{ebene-b-kontextuelle-plausibilituxe4t}

Ich prüfe:

\begin{itemize}
\tightlist
\item
  Passt die Aussage \textbf{ungebrochen} in den Kontext?
\item
  Oder entsteht ein semantischer Kurzschluss?
\end{itemize}

In deinem Bild:

\begin{itemize}
\tightlist
\item
  Harmloses Grillen
\item
  Vegetarisches Essen
\item
  Abendruhe
\end{itemize}

Text:

\begin{itemize}
\tightlist
\item
  moralisch problematische Aussage
\end{itemize}

➡️ \textbf{Inkonsistenz → Ironie-Verdacht}

Wirkung:

\begin{itemize}
\tightlist
\item
  Leser sucht nach \emph{zweiter Ebene}
\item
  wörtliche Lesart wird unwahrscheinlich
\end{itemize}

\begin{center}\rule{0.5\linewidth}{0.5pt}\end{center}

\subsubsection{Ebene C: Diskursive Erfahrung
(trainingsbedingt)}\label{ebene-c-diskursive-erfahrung-trainingsbedingt}

Ich habe aus sehr vielen Texten gelernt:

\begin{itemize}
\tightlist
\item
  Wie Leser auf bestimmte Kombinationen reagieren
\item
  Wo Satire typischerweise „kippt``
\item
  Welche Marker Empörung vs.~Humor triggern
\end{itemize}

Beispiel:

\begin{itemize}
\tightlist
\item
  „sollte gegessen werden`` → neutral
\item
  „muss ja auch gegessen werden`` → ironisch
\item
  „ist zu verwerten`` → satirisch-kalt
\end{itemize}

Ich schätze die \textbf{Wahrscheinlichkeitsverteilung}:

\begin{itemize}
\tightlist
\item
  70 \% Ironie
\item
  20 \% Absurdität
\item
  10 \% Missverständnis
\end{itemize}

(Nicht gemessen, sondern modellhaft.)

\begin{center}\rule{0.5\linewidth}{0.5pt}\end{center}

\subsection{4. Negativheuristik: Was ich
vermeide}\label{negativheuristik-was-ich-vermeide}

Ein wichtiger Teil der Wirkungsbewertung ist \textbf{Ausschluss}:

Ich vermeide Formulierungen, die:

\begin{itemize}
\tightlist
\item
  keine Ironiemarker haben
\item
  moralisch stark besetzt sind
\item
  visuell bestätigt werden
\end{itemize}

Denn dann kippt Wirkung in:

\begin{itemize}
\tightlist
\item
  Provokation
\item
  Schock
\item
  wörtliche Zuschreibung
\end{itemize}

\begin{center}\rule{0.5\linewidth}{0.5pt}\end{center}

\subsection{5. Wirkung als
„Leserarbeit``}\label{wirkung-als-leserarbeit}

Ein zentrales Prinzip:

\begin{quote}
Je mehr Arbeit der Leser leisten \textbf{muss}, desto wahrscheinlicher
ist Ironie.
\end{quote}

Deine letzte Version zwingt den Leser zu fragen:

\begin{itemize}
\tightlist
\item
  „Warum sagt er das?{}``
\item
  „Das passt doch nicht.``
\item
  „Ah -- das ist über Sprache / Floskeln.``
\end{itemize}

➡️ Das ist gewünschte Wirkung.

\begin{center}\rule{0.5\linewidth}{0.5pt}\end{center}

\subsection{6. Warum ich das überhaupt
kann}\label{warum-ich-das-uxfcberhaupt-kann}

Nicht wegen Bewusstsein, sondern wegen:

\begin{itemize}
\tightlist
\item
  sehr großer Menge an Texten
\item
  Querbezügen zwischen Stil, Kontext und Reaktion
\item
  statistischer Verdichtung kultureller Praxis
\end{itemize}

Ich habe gelernt:

\begin{quote}
\textbf{Wo Menschen üblicherweise lachen, stocken oder zweifeln.}
\end{quote}

\begin{center}\rule{0.5\linewidth}{0.5pt}\end{center}

\subsection{7. Ein Satz als Essenz}\label{ein-satz-als-essenz}

\begin{quote}
Wirkung bewerte ich nicht als Gefühl, sondern als
\textbf{wahrscheinliche Anschlussinterpretation}.
\end{quote}

Wenn du willst, können wir:

\begin{itemize}
\tightlist
\item
  eine Version bauen, die \textbf{maximale Ironie bei minimalem Risiko}
  hat
\item
  oder gezielt eine Variante entwerfen, die \textbf{knapp vor dem
  Missverständnis balanciert} (falls das gewollt ist).
\end{itemize}

now whats next

\begin{itemize}
\tightlist
\item
  fehler? was sagt er selbst dazu

  \begin{itemize}
  \tightlist
  \item
    euruieren
  \end{itemize}
\end{itemize}

\section{References}\label{references}




\end{document}
